% !TEX root = ../main.tex
\section{Empirical Results of Prompting, Evaluation, and Failure Modes}
\label{sec:prompting-failures}
\label{sec:root-cause}

This section answers RQ3 and RQ4. We first measure whether richer task specifications help, then examine evaluator reliability and the structure of non-SE-PASS failures.

\subsection{Effect of Prompt Granularity (RQ3)}
\label{sec:results-prompt}

Each task was run with a short ticket-style prompt and a long engineering-proposal (EP) prompt. Table~\ref{tab:prompt} aggregates the effect by agent. No short-prompt run reaches SE-PASS. Long prompts improve mean composite scores for all four agents, but the SE-PASS increase comes only from bounded calibration tasks and T10.

\begin{table*}[t]
\caption{Prompt effect. $\Delta$ = long $-$ short, percentage points for SE-PASS and score points for mean.}
\label{tab:prompt}
\centering
\small
\begin{tabular*}{\textwidth}{@{\extracolsep{\fill}}lrrrrrr@{}}
\toprule
\textbf{Agent} & \textbf{Short SE-PASS} & \textbf{Long SE-PASS} & \textbf{$\Delta$SE-PASS} & \textbf{Short mean} & \textbf{Long mean} & \textbf{$\Delta$mean} \\
\midrule
Claude   & 0.0\% & 4.8\% & +4.8 & 5.08 & 6.28 & +1.20 \\
Codex    & 0.0\% & 4.8\% & +4.8 & 5.20 & 6.44 & +1.24 \\
Cursor   & 0.0\% & 1.6\% & +1.6 & 4.00 & 5.77 & +1.77 \\
OpenCode & 0.0\% & 1.6\% & +1.6 & 3.72 & 4.60 & +0.88 \\
\bottomrule
\end{tabular*}
\end{table*}

On the main corpus, all four agents move from 0/50 short-prompt SE-PASS to 1/50 long-prompt SE-PASS, but the single passing task is always T10. Excluding T10, long prompts improve mean score by +0.99 to +2.12 across agents while producing zero SE-PASS. Thus, richer specifications improve partial draft quality but do not overcome the coordination barrier in non-trivial long-horizon tasks.

This result should not be read as proving that the short prompt lacked all necessary information. The EP prompt adds motivation, constraints, affected subsystems, non-goals, and expected checks, but it still avoids leaking GT file paths or implementation structure. The remaining gap is therefore consistent with two explanations that our design cannot fully separate: agents may still lack task-critical repository information, and they may fail to operationalize a long specification into a correct cross-module plan. The narrower conclusion is that structured prompt detail improves drafts, but prompt detail alone does not close the system-level feature gap.

\finding{3}{Long prompts improve drafts, but do not resolve whether failures come from missing task information or failed long-context operationalization. They raise mean scores for every agent, yet produce no merge-ready output on the 49 non-trivially scoped main-corpus tasks.}

\subsection{Reliability of LLM-Based Evaluation (RQ4)}
\label{sec:results-agreement}

Pairwise verdict agreement among the four LLM evaluators ranges from 42.5\% (Codex--GLM) to 74.4\% (Claude--Cursor); full four-way consensus holds for 33.7\% of runs. Because single-judge disagreement is substantial, all headline verdicts use the four-judge majority rule described in \S\ref{sec:eval-framework}.

\par\noindent\fbox{\parbox{0.96\columnwidth}{\small%
\textbf{Human calibration summary.}
The 12-task calibration subset contains 96 runs, each blind-rated by three senior reviewers with domain expertise.
SE-PASS-vs-non-SE-PASS agreement between the human reviewers and the LLM panel is \textbf{92.7\%} ($\kappa{=}0.43$).
The LLM panel is stricter: it never promotes a human-rejected run to SE-PASS and demotes 6 of 9 human-SE-PASS runs.}}\par

We also test same-family bias by recomputing each agent's headline verdict after removing the judge from the same model family. SE-PASS counts are unchanged for Claude (3$\to$3), Codex (3$\to$3), and Cursor (1$\to$1); OpenCode drops from 1 to 0. Only 6.9\% of verdict cells change under leave-one-family-out, and Codex's own judge is the strictest evaluator of Codex output. These checks do not remove all LLM-as-judge threats, but they indicate that the 1.0\% main-corpus SE-PASS rate is not caused by same-family judge leniency.

\finding{4}{A single LLM judge is too unstable for open-ended feature-PR evaluation, but a conservative multi-judge panel calibrated against human reviewers gives a reproducible lower-bound estimate of merge readiness.}

\subsection{Failure Signatures (RQ4)}
\label{sec:failure-modes}

For each non-SE-PASS run, we average the four judges' A/B/C scores and classify the run using the first matching rule in Table~\ref{tab:failure-signatures}. The classification is computed from the per-run evaluator CSV used for the headline tables.

\begin{table}[t]
\centering
\scriptsize
\caption{Failure signatures over the 396 non-SE-PASS main-corpus runs.}
\label{tab:failure-signatures}
\begin{tabularx}{\columnwidth}{@{}lYrr@{}}
\toprule
\textbf{Signature} & \textbf{Rule} & \textbf{Runs} & \textbf{Share} \\
\midrule
Semantic failure & $\bar{A} < 2$ & 177 & 44.7\% \\
Coverage gap & $\bar{A} \ge 2, \bar{B} < 2$ & 101 & 25.5\% \\
Design mismatch & $\bar{A},\bar{B} \ge 2, \bar{C} < 2$ & 43 & 10.9\% \\
Borderline partial & all $\ge 2$, not SE-PASS & 75 & 18.9\% \\
\bottomrule
\end{tabularx}
\end{table}

The largest bucket is semantic failure, so the benchmark is not simply exposing incomplete file discovery. At the same time, coverage gaps and borderline partials account for 44.4\% of failures, showing that a meaningful minority of non-SE-PASS outputs contain recoverable structure. Table~\ref{tab:failure-agent} further shows that agents fail differently: Codex has the largest borderline bucket, while OpenCode has the largest semantic-failure bucket.

\begin{table}[t]
\centering
\scriptsize
\caption{Failure signatures by agent on the main corpus, excluding SE-PASS runs ($n{=}99$ per agent).}
\label{tab:failure-agent}
\begin{tabular*}{\columnwidth}{@{\extracolsep{\fill}}lrrrr@{}}
\toprule
\textbf{Agent} & \textbf{Semantic} & \textbf{Coverage} & \textbf{Borderline} & \textbf{Design} \\
\midrule
Claude   & 34 & 25 & 21 & 19 \\
Codex    & 29 & 31 & 30 & 9 \\
Cursor   & 45 & 27 & 21 & 6 \\
OpenCode & 69 & 18 & 3  & 9 \\
\bottomrule
\end{tabular*}
\end{table}

\finding{5}{Non-SE-PASS outputs are dominated by semantic failures (44.7\%), but coverage gaps (25.5\%) and borderline partials (18.9\%) identify a substantial recoverable minority. This supports scope/test validation as a practical filter, not as a replacement for stronger agent reasoning.}

\subsection{Engineering-Addressable and Capability-Limited Failures}
\label{sec:addressable-failures}

The failure taxonomy separates two kinds of improvement. Coverage gaps and some borderline partials are \emph{engineering-addressable}: a better harness can require the agent to enumerate affected files, generate or update tests, run static checks, and explain why the implementation reaches the feature boundary. These checks do not make the agent smarter, but they can prevent early stopping from being mistaken for a reviewable PR. In the main corpus, this addressable region covers at least the 101 coverage-gap failures and part of the 75 borderline partials.

Semantic failures are different. When $\bar{A}<2$, adding more file-discovery scaffolding is unlikely to repair the run because the agent has misunderstood the feature, the invariant, or the implementation layer. These failures motivate model and agent-design improvements: stronger repository-level semantic reasoning, better plans that preserve invariants across edits, and validation loops that test behavior rather than only syntax. This distinction is why the discussion frames scope/test gates as a burden-reduction mechanism rather than a path to autonomous system-level feature delivery.
